\documentclass[11pt]{article}

\usepackage[margin=0.5cm]{geometry}
\usepackage{amsmath}
\usepackage{amssymb}
\usepackage{array}
\usepackage{enumitem}
\usepackage{graphicx}

\setlength{\columnsep}{0.55cm}
\setlength{\parindent}{0pt}

\newcommand{\psfig}[2]{%
  \begin{center}
  \includegraphics[width=#2]{#1}
  \end{center}
}

\begin{document}

\twocolumn

\title{Problem Set 5, COSC330/PHYS306\\Chapters 6--9}
\author{Prof. Jordan C. Hanson}
\date{}
\maketitle

\section{Chapter 6}

\begin{enumerate}

\item \textbf{[Section 6-4, Exercise 13]}\\
The waveforms in Figure 6--73 are applied to the comparator as shown.
Determine the $(A=B)$ waveform.

\psfig{fig_6_73.png}{0.98\linewidth}
\begin{center}\textbf{FIGURE 6--73}\end{center}

\item \textbf{[Section 6-4, Exercise 14]}\\
For the 4-bit comparator in Figure 6--74, plot each output waveform for
the inputs shown. The outputs are active-HIGH.

\psfig{fig_6_74.png}{0.98\linewidth}
\begin{center}\textbf{FIGURE 6--74}\end{center}

\item \textbf{[Section 6-5, Exercise 20]}\\
If the input waveforms are applied to the decoding logic as indicated in
Figure 6--76, sketch the output waveform in proper relation to the inputs.

\psfig{fig_6_76.png}{0.98\linewidth}
\begin{center}\textbf{FIGURE 6--76}\end{center}

\item \textbf{[Section 6-8, Exercise 28]}\\
For the multiplexer in Figure 6--79, determine the output for the following
input states:
\begin{equation*}
D_0=0,\quad D_1=1,\quad D_2=1,\quad D_3=0,
\end{equation*}
\begin{equation*}
S_0=1,\quad S_1=0.
\end{equation*}

\psfig{fig_6_79.png}{0.78\linewidth}
\begin{center}\textbf{FIGURE 6--79}\end{center}

\item \textbf{[Section 6-8, Exercise 30]}\\
The waveforms in Figure 6--81 are observed on the inputs of a 74HC151
8-input multiplexer. Sketch the $Y$ output waveform.

\psfig{fig_6_81.png}{0.91\linewidth}
\begin{center}\textbf{FIGURE 6--81}\end{center}

\item \textbf{[Special Design Problems, Exercise 46]}\\
Implement the logic function specified in Table 6--14 by using a 74HC151
data selector.

\begin{center}
\renewcommand{\arraystretch}{1.05}
\begin{tabular}{c c c c | c}
\hline
\multicolumn{4}{c|}{\textbf{Inputs}} & \textbf{Output} \\
$A_3$ & $A_2$ & $A_1$ & $A_0$ & $Y$ \\
\hline
0&0&0&0&0\\
0&0&0&1&0\\
0&0&1&0&1\\
0&0&1&1&1\\
0&1&0&0&0\\
0&1&0&1&0\\
0&1&1&0&1\\
0&1&1&1&1\\
1&0&0&0&1\\
1&0&0&1&0\\
1&0&1&0&1\\
1&0&1&1&1\\
1&1&0&0&0\\
1&1&0&1&1\\
1&1&1&0&0\\
1&1&1&1&1\\
\hline
\end{tabular}

\vspace{0.10cm}
\textbf{TABLE 6--14}
\end{center}

\end{enumerate}

\section{Chapter 7}

\begin{enumerate}
\setcounter{enumi}{6}

\item \textbf{[Section 7-1, Exercise 7]}\\
For a gated D latch, the waveforms shown in Figure 7--76 are observed on
its inputs. Draw the timing diagram showing the output waveform you would
expect to see at $Q$ if the latch is initially RESET.

\psfig{fig_7_76.png}{0.96\linewidth}
\begin{center}\textbf{FIGURE 7--76}\end{center}

\item \textbf{[Section 7-2, Exercise 8]}\\
Two edge-triggered J-K flip-flops are shown in Figure 7--77. If the inputs
are as shown, draw the $Q$ output of each flip-flop relative to the clock,
and explain the difference between the two. The flip-flops are initially
RESET.

\psfig{fig_7_77.png}{0.98\linewidth}
\begin{center}\textbf{FIGURE 7--77}\end{center}

\item \textbf{[Section 7-2, Exercise 14]}\\
Determine the $Q$ waveform relative to the clock if the signals shown in
Figure 7--83 are applied to the inputs of the J-K flip-flop. Assume that
$Q$ is initially LOW.

\psfig{fig_7_83.png}{0.98\linewidth}
\begin{center}\textbf{FIGURE 7--83}\end{center}

\item \textbf{[Section 7-7, Exercise 33]}\\
A 74HC00 quad NAND gate IC is used to construct a gated S-R latch on a
protoboard in the lab as shown in Figure 7--93. The schematic in part (a)
is used to connect the circuit in part (b). When you try to operate the
latch, you find that the $Q$ output stays HIGH no matter what the inputs
are. Determine the problem.

\psfig{fig_7_93.png}{0.98\linewidth}
\begin{center}\textbf{FIGURE 7--93}\end{center}

\end{enumerate}

\section{Chapter 8}

\begin{enumerate}
\setcounter{enumi}{10}

\item \textbf{[Section 8-2, Exercise 14]}\\
The shift register in Figure 8--54(a) has \(\overline{\mathrm{SHIFT/LOAD}}\)
and CLK inputs as shown in part (b). The serial data input (SER) is a 0.
The parallel data inputs are $D_0=1$, $D_1=0$, $D_2=1$, and $D_3=0$ as
shown. Develop the data-output waveform in relation to the inputs.

\psfig{fig_8_54.png}{0.98\linewidth}
\begin{center}\textbf{FIGURE 8--54}\end{center}

\end{enumerate}

\section{Chapter 9}

\begin{enumerate}
\setcounter{enumi}{11}

\item \textbf{[Section 9-2, Exercise 3]}\\
For the ripple counter shown in Figure 9--65, show the complete timing
diagram for eight clock pulses, showing the clock, $Q_0$, and $Q_1$
waveforms.

\psfig{fig_9_65.png}{0.98\linewidth}
\begin{center}\textbf{FIGURE 9--65}\end{center}

\end{enumerate}

\end{document}
